\chapterhead{Testing and Performance}

Testing ran at three levels. Unit tests checked the decision logic on its own.
Integration tests checked that the website and the server agree with each other.
System tests exercised the assembled hardware end to end, from a finger at the
door to a borrowed item appearing in the record.

One limit applies to everything below. The tests establish that each behaviour
happens, not how fast or how reliably it happens over many attempts. Timing and
error rates were not recorded across a controlled set of trials, so the rows
marked \textbf{[MEASURE]} in Table~\ref{tab:results} and
Table~\ref{tab:timing} are left for that work.

\sectionhead{Unit and Integration Testing}

Unit testing means checking small pieces of code one at a time. We wrote 18 of
these tests for the server, where the decisions are made. They cover three
things: that passwords follow the rules, that the warning checks work on a quiet
record as well as a busy one, and that the date maths is right.

Integration testing means checking that two parts agree with each other. The
website is checked against the server automatically. That check passed with no
errors across 4{,}871 files. This catches the case where the server starts
sending something the website does not expect.

The reader firmware was compiled for the target board before flashing. Both
programs built with no errors, which confirmed that the modules link correctly
and that the libraries match the board.

\sectionhead{System Testing}

System testing means using the whole thing the way a person would. We ran the
server and the website together against a database that already held records.
This matters. An empty database returns nothing, so a fault that only appears
with real data would stay hidden.

We logged in first, on purpose. It worked and returned the right account. This
proved the system was alive, so any later failure would mean something real and
not just a dead server.

Then we opened each page the website uses. The record list, the borrow list, the
active session check, and the live feed all worked. The record list held 39
saved events of eight different kinds. All 20 pages of the website loaded.

\subsectionhead{Testing the Assembled Readers}

With both readers built and flashed, the complete sequence was exercised on the
hardware itself. A registered finger at the door released the door lock and
opened a session. The same finger at the cabinet released the cabinet within
that session. The cabinet displayed its one-time code, and the website accepted
a borrow record only when that code was entered.

The refusal paths were exercised as well, since a system that only ever grants
access has not been tested. A registered finger presented at the cabinet without
a door session first was refused, and the refusal was written to the record as a
tailgating attempt. A finger that had expired out of its session was likewise
refused. In both cases the cabinet stayed shut and the attempt appeared in the
trail with a name, a reader, and a time.

\sectionhead{Test Results}

Table~\ref{tab:results} records the outcome of every activity described above,
together with the evidence behind it. Rows marked \textbf{[MEASURE]} are
behaviours already shown to work whose performance has not yet been quantified.

\begin{table}[H]
\centering
\caption{What we tested and what happened}
\label{tab:results}
% APA permits single spacing within a table. Needed here: at double
% spacing this one runs a few points past the bottom of the page.
\linespread{1}\selectfont
\begin{tabular}{@{}p{46mm}p{26mm}p{42mm}@{}}
\toprule
What we tested & Result & How we know \\
\midrule
\multicolumn{3}{@{}l}{\textit{Software}} \\
Server unit tests      & Passed & 18 tests, none failed \\
Website and server fit & Passed & 0 errors, 4{,}871 files \\
Door reader build      & Passed & Compiled, no errors \\
Cabinet reader build   & Passed & Compiled, no errors \\
Logging in             & Passed & Correct account returned \\
Record list            & Passed & 39 events, 8 kinds \\
Borrow list            & Passed & Records returned \\
Session check          & Passed & Worked \\
Live feed              & Passed & Stayed connected \\
Website pages          & Passed & 20 of 20 loaded \\
\midrule
\multicolumn{3}{@{}l}{\textit{Assembled hardware}} \\
Door scan opens session  & Passed & Lock released, session created \\
Cabinet opens for same person & Passed & Released inside the session \\
One-time code shown      & Passed & Displayed on opening \\
Borrow record needs code & Passed & Refused without it \\
Tailgating refused       & Passed & No session, cabinet stayed shut \\
Expired session refused  & Passed & Cabinet stayed shut \\
Refusals recorded        & Passed & Written to the trail \\
\midrule
\multicolumn{3}{@{}l}{\textit{Not yet quantified}} \\
Scan time            & [MEASURE] & Works; not timed over trials \\
Unlock time          & [MEASURE] & Works; not timed over trials \\
False acceptance rate & [MEASURE] & Needs a controlled trial set \\
False rejection rate  & [MEASURE] & Needs a controlled trial set \\
Text messages        & Stand-in only & No real SIM card \\
\bottomrule
\end{tabular}
\end{table}

We found one fault, and we report it instead of quietly fixing it, because it
shows why testing with real data matters. The server sent all 39 records, but
the website showed none of them. The reason: the server had learned two new
kinds of record that the website did not know about. The website checks the
whole list at once, so one unknown item threw away all 39. The spreadsheet page,
which checks each row on its own, showed everything. That difference is what
told us where the fault was. The saved records themselves were never harmed.

\sectionhead{System Requirements}

Table~\ref{tab:reqs} states the requirements the system was designed against,
separated into what it must do and how it must behave. These are the statements
the tests above were written to check.

\begin{table}[H]
\centering
\caption{What the system must do}
\label{tab:reqs}
\linespread{1}\selectfont
\begin{tabular}{@{}p{10mm}p{100mm}@{}}
\toprule
No. & Requirement \\
\midrule
\multicolumn{2}{@{}l}{\textit{What it must do}} \\
F1 & The door opens only for a finger signed up on the door sensor. \\
F2 & A door scan starts a session for that one person, which ends by itself. \\
F3 & The cabinet opens only for the person holding a running session. \\
F4 & Every decision is saved with the person, the reader, and the time. \\
F5 & The cabinet shows a one-time code, and no borrow record is accepted
     without it. \\
F6 & The website shows which readers are working. \\
\midrule
\multicolumn{2}{@{}l}{\textit{How it must behave}} \\
N1 & Readers never decide on their own. \\
N2 & Fingerprints never leave the sensor. \\
N3 & Nobody can edit or delete a saved reader record, not even an
     administrator. \\
N4 & Each reader has its own password, which can be changed on its own. \\
N5 & The program must fit in the chip with room to spare. \\
N6 & The reader must never freeze while the lock is open or a message is
     showing. \\
\bottomrule
\end{tabular}
\end{table}

Every requirement in the table was exercised during testing. F1 to F6 were
demonstrated on the assembled readers, as recorded in
Table~\ref{tab:results}. N1, N2 and N4 follow from how the system is built and
were confirmed by inspection: a reader holds no decision logic, no template
leaves the sensor, and each reader carries its own credentials. N3 was confirmed
by attempting to edit a reader record through the interface, which offers no
route to do so. N5 is measured below. N6 was observed on the hardware, where a
second finger could be presented while a message was still on screen.

\sectionhead{Performance Analysis}

Performance is examined here in two ways. The first is what the finished program
demands of the chip, which the build reports exactly. The second is how the
design sits against systems that have already been measured and published.

\subsectionhead{Use of Chip Resources}

Our program fills about 81 percent of its 1.2 MB space. That meets requirement
N5, but only just. Adding features later will be limited by this, not by speed
or memory.

RAM is not a problem, and the reason is a design choice. Fingerprints stay in
the sensor, so the biggest piece of data never reaches our chip. That turns the
usual tightest limit on this kind of project into a non-issue. It also answers
the small-memory warning raised by \textcite{rahman2026}.

Speed is not a problem either. \textcite{espressif2025} rate the chip at 1079.96
CoreMark, a standard speed score. Our reader only waits on a slow sensor, builds
a short message, and keeps Wi-Fi alive. Two processors mean the Wi-Fi keeps
running while the sensor is read.

\subsectionhead{Timing}

Table~\ref{tab:timing} separates the intervals fixed by the design, which are
known exactly, from those that depend on the sensor and the network and would
have to be timed over repeated trials.

\begin{table}[H]
\centering
\caption{Timings we set, and timings still to be measured}
\label{tab:timing}
\begin{tabular}{@{}p{42mm}p{26mm}p{42mm}@{}}
\toprule
Interval & Value & Status \\
\midrule
Session length        & 120~s     & Set in the server \\
Checkout code lifetime & 900~s    & Set in the server \\
Enrollment code lifetime & 300~s  & Set in the server \\
Door lock stays open  & 3~s       & Set in the reader \\
Cabinet lock stays open & 4~s     & Set in the reader \\
Message on screen     & 2.5~s     & Set in the reader \\
``Still here'' signal & 30~s      & Set in the reader \\
Scan time             & [MEASURE] & Observed to work; not timed \\
Unlock time           & [MEASURE] & Observed to work; not timed \\
\bottomrule
\end{tabular}
\end{table}

\subsectionhead{Comparison With Published Systems}

The figures below belong to other people's systems. They are the benchmarks
this system should be compared against once its own numbers are recorded.

\textcite{alsayaydeh2025} measured a scan at about one second and an unlock at
five seconds on a setup much like ours. They admitted the wrong person 1.32
percent of the time and turned away the right person 26.32 percent of the time.
The second figure is the warning. A system that rejects one honest person in
four will end up propped open, and a propped door defeats everything behind it.

\textcite{mohammed2025} report a mean authentication time of 235 milliseconds
and a 68 percent drop in fraudulent access. Their timing is not a fair
comparison, since their pipeline runs on server-class machines rather than a
microcontroller. The 68 percent is the relevant figure: it is the scale of
improvement that checking twice is expected to deliver, and the effect the
door-to-cabinet sequence is built to reproduce.

Two properties of this system have no counterpart in either study. Refusals are
recorded as carefully as approvals, so tailgating attempts become a countable
quantity rather than an assumption, and the refusal tests above confirm those
records are written. Reader events are also append-only by construction, which
is the accountability property \textcite{koisser2023} argue a record must have
before it can be used as evidence.
